\documentclass[11pt]{article}
\usepackage[margin=1in]{geometry}
\usepackage{lmodern}
\usepackage[T1]{fontenc}
\usepackage[utf8]{inputenc}
\usepackage{microtype}
\usepackage{xcolor}
\usepackage{hyperref}
\hypersetup{
  colorlinks=true,
  linkcolor=black,
  urlcolor=blue!55!black,
  citecolor=blue!55!black,
}
\usepackage{booktabs}
\usepackage{array}
\usepackage{listings}
\lstset{
  basicstyle=\ttfamily\footnotesize,
  breaklines=true,
  columns=fullflexible,
  backgroundcolor=\color{gray!7},
  frame=none,
  xleftmargin=1em,
  aboveskip=0.6em,
  belowskip=0.6em,
}
\usepackage{enumitem}
\setlist{topsep=2pt, itemsep=1pt}
\setlength{\parskip}{0.55em}
\setlength{\parindent}{0pt}
\usepackage{titlesec}
\titlespacing*{\section}{0pt}{1.4em}{0.5em}
\titlespacing*{\subsection}{0pt}{1.0em}{0.3em}

\title{\textbf{\LARGE Pothi: owning your bibliography} \\[0.5em]
       \large \textit{Your reference manager belongs in your folder, \\
                      not in someone else's cloud.}}
\author{Kanchan Sarkar \\[0.3em]
        \small \texttt{pcksarkar@gmail.com}}
\date{\today}

\begin{document}
\maketitle

\begin{abstract}
Most reference managers --- Zotero, Mendeley, EndNote --- treat your library as
one big pile and lock it inside their software. \textbf{Pothi} is built
differently: each paper you are \emph{writing} is its own folder on disk, and
that folder owns its own auto-updating \texttt{references.bib}. Drop a PDF
into the folder and Pothi finds it, looks up the metadata from CrossRef, and
adds the entry to your library. \textbf{If you write in LaTeX}, Pothi keeps
\texttt{references.bib} in your manuscript's folder current as you cite ---
your next \texttt{pdflatex} run picks up the change with no manual export
step. \textbf{If you write in Word}, type \texttt{[@SmithCrystal2024]} in
your document, drop the docx into Pothi, and out comes a polished version
with formatted citations and a bibliography. No account to make, no server
to run, no subscription --- your library is plain JSON sitting in a folder
you chose. This document explains the design choices behind the tool, walks
through five things you would actually do with it (add a paper by DOI, drag
in a PDF, build a manuscript bibliography for LaTeX or Word, recover after
switching browsers), compares it honestly to the popular options, and
includes the installation and usage manual. The code is open under the MIT
license. I wrote it for myself because the existing tools did not fit how I
work --- sharing it because it might fit you too.
\end{abstract}

\section{Why I built this}

I write papers in LaTeX, sometimes Word. My references live in PDFs
scattered across several folders on my laptop and my Nextcloud. I tried
Zotero and Mendeley over the years; both eventually frustrated me for
the same reasons.

\begin{itemize}
\item \textbf{Cloud lock-in.} Both want my data in their cloud. Mendeley
sold to Elsevier; Zotero is community-run but its sync is server-side. I
have lost a Mendeley library before to a sync conflict.
\item \textbf{Library-first model.} Both organize the world as one global
library that I then ``use'' to cite from. But when I am writing a paper, the
unit of work is \emph{the paper}, not the library. I want this paper's
bibliography to be a curated, ordered list with my own notes about why each
reference is in --- and I want that list to be a plain file next to my
manuscript, not a tag inside a database.
\item \textbf{Attachment fragility.} Rename a PDF on disk and the link in
Zotero breaks. Move a folder and dozens of links go stale.
\item \textbf{Vendor pace.} EndNote costs money. Papers and ReadCube changed
hands. JabRef is great if you write in LaTeX; awkward otherwise.
\item \textbf{The \texttt{.bib} file should live next to my manuscript.}
When I write a LaTeX paper, the natural unit is the folder I
\texttt{pdflatex} from. I want \texttt{references.bib} to sit in that
folder, kept current as I cite, and to be the same file I email a
co-author. Most tools instead want a global library that I ``cite from''
through a plugin.
\end{itemize}

I am a researcher, not a software person. So this is not a polished commercial
tool; it is a thing I wrote for myself in a few hundred hours over a few
weeks. It runs entirely in the browser, has no backend, and stores everything
in a folder you pick. If it fits how you work, take it. If not, the existing
tools are right there.

\section{Design principles}

Five principles kept the project small and consistent.

\subsection{Local-first}

The data is yours. There is no account, no server, no subscription. The
library is plain JSON in a folder you control. If Pothi disappears, your
files are still there. Inspired by the Local-First Software essay (Kleppmann
et al., Ink \& Switch, 2019).

\subsection{Manuscript-first}

A reference is interesting because of what \emph{paper} you cite it from.
Pothi treats each manuscript as a first-class object: it has a name, an
associated folder, an ordered list of cited references, and a per-citation
rationale (``why is this cited?''). The global library is the union of all
manuscript bibliographies, not the other way around.

\subsection{Plain text on disk}

When you link a folder, Pothi auto-writes \verb|_pothi-library.json| (the
whole library) and per-manuscript \verb|references.bib| / \verb|refs.json|
files into it. Open them in any text editor. Track them in git. Email them.
The browser's IndexedDB is a cache, not the source of truth.

\subsection{Schema as config}

Entry types, required fields, citekey templates --- they all live in
\texttt{js/schema.js} as a plain configuration object. To add a custom entry
type or field, edit the file. There is no plugin API to learn.

\subsection{No build step}

The project is HTML, CSS, and ES module JavaScript. Edit a file, refresh the
browser. Vendored libraries (Preact, htm, PDF.js, fflate) live in
\texttt{vendor/} as committed modules. There is no \texttt{npm install}, no
webpack, no transpilation.

\section{How it is put together}

The full project is one folder of static files.

\begin{lstlisting}
Pothi/
  index.html              # shell + import map
  css/
    variables.css         # design tokens
    base.css, layout.css, components.css
  js/
    schema.js             # entry types + fields + citekey template
    db.js                 # IndexedDB wrapper (entries, manuscripts, meta)
    bibtex.js             # parser + emitter, with LaTeX -> Unicode
    citekey.js            # {Author1}{year}{title3} template engine
    lookup.js             # CrossRef + arXiv (via DOI) + OpenLibrary + Sem.Sch.
    folder.js             # File System Access API wrapper
    folder-scan.js        # recursive walk + SHA-256 dedup
    pdf-extract.js        # PDF.js + DOI/arXiv regex
    pdf-preview.js        # in-app PDF viewer
    docx.js               # Pandoc-bridge equivalent (fflate + DOMParser)
    styles.js             # author-year + numeric citation styles
    library-backup.js     # export/import + auto-write to folder
    smart-collections.js  # saved searches
    manuscripts.js        # manuscript helpers
    markdown.js           # tiny safe Markdown -> DOM
    tag-suggest.js        # TF-IDF tag suggestions
    app.js                # main render
  vendor/                 # Preact, htm, PDF.js, fflate (vendored)
  install.sh, install.bat # launcher installers
  LICENSE                 # MIT
\end{lstlisting}

The browser entry point is one HTML file with an \emph{import map} that points
each module's name at its vendored file:

\begin{lstlisting}
<script type="importmap">
{ "imports": {
    "preact":       "./vendor/preact.module.js",
    "preact/hooks": "./vendor/preact-hooks.module.js",
    "htm":          "./vendor/htm.module.js"
  } }
</script>
\end{lstlisting}

IndexedDB stores three object stores: \texttt{entries}, \texttt{manuscripts},
and a \texttt{meta} key/value singleton (where the linked-folder
\texttt{FileSystemDirectoryHandle} lives across reloads).

The browser refuses ES modules from \texttt{file://} for security, so Pothi
needs an HTTP server. It ships an installer (\texttt{install.sh} on Linux/
macOS, \texttt{install.bat} on Windows) that drops a \texttt{pothi} command
on your \texttt{PATH}: typing \texttt{pothi} in any terminal starts a Python
\texttt{http.server} on \texttt{127.0.0.1:8765} and opens your browser there.

\section{Five things you would actually do with it}

\subsection{Add a paper by DOI / arXiv ID / ISBN / URL}

Click \emph{+ New}. Paste any of: a bare DOI (\texttt{10.1038/nature14539}),
a \texttt{doi:} prefix, a full URL (\texttt{https://doi.org/10.x/y} or any
journal page that has a DOI in it --- the input sniffs out the DOI), an arXiv
ID (\texttt{2401.04088}), or an ISBN (\texttt{9780198506058}). Press Enter.
The handler routes to CrossRef / arXiv-via-CrossRef / OpenLibrary as
appropriate. Title, authors (in BibTeX \emph{Last, First} form), journal,
volume, issue, pages, year, DOI all populate. When CrossRef has no abstract,
Pothi falls back to Semantic Scholar to fill it. A citekey is generated from
\texttt{\{Author1\}\{year\}\{title3\}} (configurable). If a duplicate DOI is
already in your library, you get a clear warning before adding.

\subsection{Drag a PDF anywhere on the page}

A blue overlay appears: ``Drop PDFs or .docx here.'' Release. Pothi:
\begin{enumerate}
\item extracts the embedded text and XMP metadata from the PDF (PDF.js);
\item regex-sniffs for a DOI on the first three pages;
\item if found, calls CrossRef to fetch full metadata;
\item generates a citekey;
\item if a library folder is linked, copies the PDF into it as
      \texttt{<citekey>.pdf} and attaches it to the new entry.
\end{enumerate}
Hit rate in practice is around 85\% for journal articles and higher for
arXiv preprints. PDFs without a DOI get the title from the PDF's embedded
metadata; you fill the rest.

If the PDF's content hash matches a file already attached to an entry, the
import is silently de-duplicated.

\subsection{Build a manuscript bibliography}

Click \emph{+ New manuscript} in the sidebar. Name it (\emph{Polymorphism
review 2026}). The main pane switches to a \emph{manuscript view}: a search
box at the top, the bibliography below.

Search the library inline; click \emph{+ Cite} to add a reference. A textarea
appears next to each cited entry where you write the rationale --- ``cited
in \S 3.2 because Bernstein establishes the canonical polymorphism
framework.'' Reorder with up / down buttons. Switch the citation style
between \emph{Author-year} (APA-ish) and \emph{Numeric} (IEEE-ish).

Click \emph{+ Link folder} on the manuscript to point it at the folder
that holds your draft (manuscripts can each have a different folder).
Pothi then auto-writes \texttt{references.bib} and \texttt{refs.json}
into that folder on every change, debounced 4 seconds.

\paragraph{If you write in LaTeX.} This is the whole workflow:

\begin{lstlisting}
% In your manuscript folder, after linking it to the manuscript:
\documentclass{article}
\usepackage{biblatex}
\addbibresource{references.bib}      % Pothi keeps this current

\begin{document}
As Bernstein argues~\cite{Bernstein2020Polymorphism}, ...
% Or for vanilla BibTeX:
% \bibliography{references}
\printbibliography
\end{document}
\end{lstlisting}

Run \texttt{pdflatex} as usual. Each time you cite or uncite in Pothi,
the \texttt{references.bib} next to your \texttt{.tex} updates within
seconds; the next compile sees the change. No export step, no Cite-While-
You-Write plugin, no app to keep running in the background. The
\texttt{.bib} file is the contract; both Pothi and your TeX engine read
and write it as plain text.

For BibLaTeX users, \texttt{references.bib} is BibLaTeX-clean (the
emitter handles Unicode, doesn't lossy-escape, preserves brace
protection). For BibTeX (8-bit) users, your usual \texttt{inputenc}
setup works.

\subsection{Write a docx with citations}

If you write in Word, type Pandoc-style placeholders directly:

\begin{itemize}
\item \texttt{[@SmithCrystal2024]} --- one citation
\item \texttt{[@Smith2024; @Doe2023]} --- multiple in one bracket
\item \texttt{[@Smith2024, p. 42]} --- with a page locator
\end{itemize}

Drag the \texttt{.docx} onto the Pothi page. Pothi unzips it
(fflate), walks every paragraph in \texttt{word/document.xml}, replaces
every \texttt{[@key]} with a formatted citation in your chosen style, and
appends a bibliography section at the end. You get a new file
\texttt{<name>-cited.docx} downloaded to your Downloads folder. Missing
citekeys show as \texttt{[?]} in the document and are listed in the
post-process toast so you can find them.

This works without installing Pandoc. If you prefer Pandoc, the export
menu also has a \emph{Use with Word + Pandoc\dots} entry that gives you
the exact command to copy-paste.

\subsection{Recover after switching browsers, ports, or machines}

Browser storage is keyed by \emph{origin}: \texttt{http://127.0.0.1:8765}
and \texttt{http://127.0.0.1:8766} are different origins, so they have
different IndexedDBs. This is the most common ``my library disappeared''
scenario.

Pothi's defense: when a library folder is linked, the auto-write of
\texttt{\_pothi-library.json} mirrors the entire library to disk on every
change (debounced). On a fresh origin, you re-link the folder; Pothi
detects the backup file, asks ``Found a library backup with N entries
(last saved $\langle$timestamp$\rangle$). Restore?'', and replays it
into the new IndexedDB. The plain JSON on disk is the source of truth;
the browser cache is rebuilt from it.

\section{An honest comparison}

\begin{table}[h]
\small
\centering
\renewcommand{\arraystretch}{1.25}
\begin{tabular}{lccccc}
\toprule
                              & Zotero+BBT & JabRef & Mendeley & EndNote & Pothi  \\
\midrule
Free / open-source            & yes / yes & yes / yes & yes / no & no / no & yes / yes \\
No account required           & yes (BBT) & yes & no  & no  & yes \\
Local-only by default         & no  & yes & no  & yes & yes \\
Plain text on disk            & no  & yes (.bib) & no  & no  & yes (JSON + .bib) \\
Manuscript-first model        & no  & no  & no  & no  & yes \\
Per-citation rationale        & no  & no  & no  & no  & yes \\
\texttt{.bib} lives next to manuscript & opt-in (BBT) & yes & no  & no  & yes \\
Live-sync \texttt{.bib} on edit & on save (BBT) & on save & no & no & yes (4\,s) \\
Watch-folder auto-import      & via ZotFile & no & no  & no  & yes \\
Drag-PDF DOI sniff            & yes & no  & yes & no  & yes \\
docx ``cite-while-write''     & yes (Word add-in) & no & yes & yes & via Pandoc bridge \\
Custom entry types            & via plugin & yes & no & yes & via schema.js \\
Runs in browser               & web only  & no  & no  & no  & yes \\
Install ceremony              & app + ext + plugin & app + Java & app + ext & app & none (browser) \\
Lines of source               & $\sim$300\,k & $\sim$200\,k & closed & closed & $\sim$10\,k \\
\bottomrule
\end{tabular}
\caption{Pothi compared to widely used reference managers. ``Yes'' / ``no''
cells are the author's reading; please correct any errors.}
\end{table}

The table is a guide, not a marketing piece. Zotero with Better BibTeX is the
right tool for most researchers; if it works for you, do not switch. Pothi
exists because a few of those columns mattered to me enough to write
something different.

\section{What it does not do}

Honesty matters here.

\begin{itemize}
\item \textbf{No Microsoft Word add-in.} The docx pipeline works through a
drop-and-replace flow, not a live task pane inside Word. Building a real
Office add-in is a separate sub-project (different deployment, Office.js,
manifest distribution) and is not in this version.
\item \textbf{Chromium-only for folder linking.} The File System Access API
is not yet in Firefox or Safari. The non-folder parts of Pothi --- adding
references manually, browsing, search, exports --- work in any modern
browser. Folder linking, watch-folder scan, and PDF attachments need
Chrome / Edge / Brave / Opera.
\item \textbf{No cloud sync.} Use Nextcloud / Dropbox / iCloud / git on the
linked folder. The plain-JSON \texttt{\_pothi-library.json} backup file is
the synchronization unit; conflicts surface as file conflicts and you
resolve them in your sync tool.
\item \textbf{No real-time multi-user.} If you open Pothi in two tabs at
the same origin they will not see each other's edits until reload. Two
machines editing the same folder concurrently will produce a sync conflict
on the backup file.
\item \textbf{Single maintainer.} I wrote this. If I get hit by a bus, the
code is MIT-licensed and small enough to fork.
\item \textbf{Citation styles are limited.} The docx pipeline ships with two
styles, author-year (APA-ish) and numeric (IEEE-ish). Bring-your-own CSL
support is a future addition; Pandoc handles it today.
\item \textbf{Semantic Scholar is rate-limited} without an API key, so the
citation-count chip is best-effort.
\end{itemize}

\section{Repository, installation, and license}

\paragraph{Code.} \url{https://github.com/<REPO>/pothi} (replace with the
real URL).

\paragraph{License.} MIT. The file \texttt{LICENSE} in the repository carries
the full text. Do anything you want with the code; just keep the copyright
line.

\paragraph{Install.} Clone the repository. From the project folder:

\begin{lstlisting}
./install.sh           # Linux / macOS
install.bat            # Windows
\end{lstlisting}

This drops a \texttt{pothi} command into your \texttt{PATH} and creates a
clickable launcher (\texttt{pothi.desktop} on Linux, \texttt{Pothi.command}
on macOS, \texttt{Pothi.bat} on Windows). After running it, reload your
shell and type \texttt{pothi}; the local server starts and your browser
opens.

\paragraph{Manual run.} If you would rather not install:

\begin{lstlisting}
python3 -m http.server 8765 --bind 127.0.0.1
\end{lstlisting}

Then open \url{http://127.0.0.1:8765/} in a Chromium-based browser.

\section{Future and open invitations}

Things on my list, in no particular order:
\begin{itemize}
\item A real Microsoft Word add-in (Office.js task pane).
\item PDF annotation extraction --- pull highlights via PDF.js, surface
them as quotable snippets attached to entries.
\item Bring-your-own CSL --- vendor citeproc-js so the docx pipeline handles
any of the 10\,000+ styles at \url{https://www.zotero.org/styles}.
\item Integration with experiment / structure databases (in my case,
crystallography records). Plain-text on disk makes this trivial in principle
but needs a target database to demonstrate.
\item A polished onboarding tour and a single-file bundled build for
double-click usage.
\end{itemize}

If you find Pothi useful, the most helpful contributions are: bug reports
with reproduction steps, screenshots of the UI breaking on your browser,
and \texttt{.bib} files where the import either dropped data or got the
Unicode wrong. Pull requests welcome.

\section*{Acknowledgements}

This project stands on the shoulders of \href{https://preactjs.com/}{Preact},
\href{https://github.com/developit/htm}{htm},
\href{https://mozilla.github.io/pdf.js/}{PDF.js},
\href{https://github.com/101arrowz/fflate}{fflate},
\href{https://www.crossref.org/}{CrossRef},
\href{https://www.semanticscholar.org/}{Semantic Scholar},
\href{https://openlibrary.org/}{OpenLibrary}, and
\href{https://pandoc.org/}{Pandoc}. The local-first framing is heavily
influenced by the Ink \& Switch
\href{https://www.inkandswitch.com/local-first/}{Local-First Software}
essay. The manuscript-first idea --- that a paper is the natural unit of
research curation, not the library --- is, as far as I know, original to
this project, but I would be surprised if no one else has independently
arrived at it.

\end{document}
